% GENERATED FILE. Edit canonical lessons, then run npm run generate:books.
% canonical-source-hash: fnv1a64:e8289e94f43b97c9
% canonical-lessons: KA-C437-erport, KA-C437-tarminal, KA-C437-bording-pas, KA-C437-skyan-madu, KA-C437-pailat

\chapter{At the Airport}
\label{ch:ka-at-the-airport}
\hlchaptermodality{437}

\begin{chapteropening}
\textbf{By the end of this chapter:} \emph{I can say an airport, a terminal, a boarding pass, to scan and a pilot.}

Complete the last lesson of chapter 437: I can say an airport, a terminal, a boarding pass, to scan and a pilot.
\end{chapteropening}

\section[\texorpdfstring{\kn{ಏರ್ಪೋರ್ಟ್} (ērpōrṭ)}{ಏರ್ಪೋರ್ಟ್ (ērpōrṭ)}]{\kn{ಏರ್ಪೋರ್ಟ್} (ērpōrṭ) — an airport}

\label{lesson:KA-C437-erport}



\begin{quote}
Before the new one: say the Kannada for hand luggage, then the Kannada for to fill up.
\end{quote}

\subsection*{You'll want to know: \kn{ಏರ್ಪೋರ್ಟ್}}
\textbf{\kn{ಏರ್ಪೋರ್ಟ್}} — \emph{ērpōrṭ} — \textquotedblleft{}an airport\textquotedblright{}.

\begin{quote}
\kn{ಏರ್ಪೋರ್ಟ್} \kn{ನಗರದಿಂದ} \kn{ತುಂಬಾ} \kn{ದೂರ} \kn{ಇದೆ}.
\end{quote}

\begin{quote}
\emph{ērpōrṭ nagaradinda tumbā dūra ide.}
\end{quote}

\begin{quote}
The airport is a long way from the city.
\end{quote}

The Kannada name is \kn{ವಿಮಾನ} \kn{ನಿಲ್ದಾಣ}, an aeroplane station; in speech the English loan is commoner.

\begin{cousinweb}
\textbf{\kn{ಏರ್ಪೋರ್ಟ್}} (\emph{ērpōrṭ}) is a loan from English \emph{airport}.
\end{cousinweb}

\begin{grammarlens}[title={What you've built}]
One more word for getting around.
\end{grammarlens}

\subsection*{Your turn}
\begin{itemize}
\raggedright
  \item \emph{Say it:} \emph{ērpōrṭ}
  \item \emph{Say it:} \emph{ērpōrṭ}, once more
  \item \emph{Say it:} \emph{ērpōrṭ nagaradinda tumbā dūra ide.} — The airport is a long way from the city.
  \item \emph{Recall:} say the Kannada for hand luggage, then the Kannada for to fill up, then say \emph{ērpōrṭ} again
\end{itemize}

\begin{tcolorbox}[breakable,colback=teal!4,colframe=teal!35!black,title={Before you move on}]
What does \kn{ಏರ್ಪೋರ್ಟ್} mean? (\textquotedblleft{}An airport\textquotedblright{}.) Say it once more.
\end{tcolorbox}
\section[\texorpdfstring{\kn{ಟರ್ಮಿನಲ್} (ṭarminal)}{ಟರ್ಮಿನಲ್ (ṭarminal)}]{\kn{ಟರ್ಮಿನಲ್} (ṭarminal) — a terminal}

\label{lesson:KA-C437-tarminal}



\begin{quote}
Before the new one: say the Kannada for to fill up, then the Kannada for an airport.
\end{quote}

\subsection*{You'll want to know: \kn{ಟರ್ಮಿನಲ್}}
\textbf{\kn{ಟರ್ಮಿನಲ್}} — \emph{ṭarminal} — \textquotedblleft{}a terminal\textquotedblright{}.

\begin{quote}
\kn{ನಮ್ಮ} \kn{ವಿಮಾನ} \kn{ಎರಡನೇ} \kn{ಟರ್ಮಿನಲ್ನಿಂದ} \kn{ಹೊರಡುತ್ತದೆ}.
\end{quote}

\begin{quote}
\emph{namma vimāna eraḍanē ṭarminalninda horaḍuttade.}
\end{quote}

\begin{quote}
Our flight leaves from the second terminal.
\end{quote}

An airport building, or the last stop of a bus line.

\begin{cousinweb}
\textbf{\kn{ಟರ್ಮಿನಲ್}} (\emph{ṭarminal}) is a loan from English \emph{terminal}, the end point of a route.
\end{cousinweb}

\begin{grammarlens}[title={What you've built}]
One more word for getting around.
\end{grammarlens}

\subsection*{Your turn}
\begin{itemize}
\raggedright
  \item \emph{Say it:} \emph{ṭarminal}
  \item \emph{Say it:} \emph{ṭarminal}, once more
  \item \emph{Say it:} \emph{namma vimāna eraḍanē ṭarminalninda horaḍuttade.} — Our flight leaves from the second terminal.
  \item \emph{Recall:} say the Kannada for to fill up, then the Kannada for an airport, then say \emph{ṭarminal} again
\end{itemize}

\begin{tcolorbox}[breakable,colback=teal!4,colframe=teal!35!black,title={Before you move on}]
What does \kn{ಟರ್ಮಿನಲ್} mean? (\textquotedblleft{}A terminal\textquotedblright{}.) Say it once more.
\end{tcolorbox}
\section[\texorpdfstring{\kn{ಬೋರ್ಡಿಂಗ್} \kn{ಪಾಸ್} (bōrḍiṅg pās)}{ಬೋರ್ಡಿಂಗ್ ಪಾಸ್ (bōrḍiṅg pās)}]{\kn{ಬೋರ್ಡಿಂಗ್} \kn{ಪಾಸ್} (bōrḍiṅg pās) — a boarding pass}

\label{lesson:KA-C437-bording-pas}



\begin{quote}
Before the new one: say the Kannada for an airport, then the Kannada for a terminal.
\end{quote}

\subsection*{You'll want to know: \kn{ಬೋರ್ಡಿಂಗ್} \kn{ಪಾಸ್}}
\textbf{\kn{ಬೋರ್ಡಿಂಗ್} \kn{ಪಾಸ್}} — \emph{bōrḍiṅg pās} — \textquotedblleft{}a boarding pass\textquotedblright{}.

\begin{quote}
\kn{ಬೋರ್ಡಿಂಗ್} \kn{ಪಾಸ್} \kn{ಇಲ್ಲದೆ} \kn{ಒಳಗೆ} \kn{ಹೋಗಲು} \kn{ಆಗುವುದಿಲ್ಲ}.
\end{quote}

\begin{quote}
\emph{bōrḍiṅg pās illade oḷage hōgalu āguvudilla.}
\end{quote}

\begin{quote}
Without a boarding pass nobody can go in.
\end{quote}

\kn{ಪಾಸ್} alone is any pass or permit card.

\begin{cousinweb}
\textbf{\kn{ಬೋರ್ಡಿಂಗ್} \kn{ಪಾಸ್}} (\emph{bōrḍiṅg pās}) is a loan from English \emph{boarding pass}, taken over whole.
\end{cousinweb}

\begin{grammarlens}[title={What you've built}]
One more word for getting around.
\end{grammarlens}

\subsection*{Your turn}
\begin{itemize}
\raggedright
  \item \emph{Say it:} \emph{bōrḍiṅg pās}
  \item \emph{Say it:} \emph{bōrḍiṅg pās}, once more
  \item \emph{Say it:} \emph{bōrḍiṅg pās illade oḷage hōgalu āguvudilla.} — Without a boarding pass nobody can go in.
  \item \emph{Recall:} say the Kannada for an airport, then the Kannada for a terminal, then say \emph{bōrḍiṅg pās} again
\end{itemize}

\begin{tcolorbox}[breakable,colback=teal!4,colframe=teal!35!black,title={Before you move on}]
What does \kn{ಬೋರ್ಡಿಂಗ್} \kn{ಪಾಸ್} mean? (\textquotedblleft{}A boarding pass\textquotedblright{}.) Say it once more.
\end{tcolorbox}
\section[\texorpdfstring{\kn{ಸ್ಕ್ಯಾನ್} \kn{ಮಾಡು} (skyān māḍu)}{ಸ್ಕ್ಯಾನ್ ಮಾಡು (skyān māḍu)}]{\kn{ಸ್ಕ್ಯಾನ್} \kn{ಮಾಡು} (skyān māḍu) — to scan}

\label{lesson:KA-C437-skyan-madu}



\begin{quote}
Before the new one: say the Kannada for a terminal, then the Kannada for a boarding pass.
\end{quote}

\subsection*{You'll want to know: \kn{ಸ್ಕ್ಯಾನ್} \kn{ಮಾಡು}}
\textbf{\kn{ಸ್ಕ್ಯಾನ್} \kn{ಮಾಡು}} — \emph{skyān māḍu} — \textquotedblleft{}to scan\textquotedblright{}.

\begin{quote}
\kn{ಅಧಿಕಾರಿ} \kn{ನನ್ನ} \kn{ಚೀಲವನ್ನು} \kn{ಸ್ಕ್ಯಾನ್} \kn{ಮಾಡಿದರು}.
\end{quote}

\begin{quote}
\emph{adhikāri nanna cīlavannu skyān māḍidaru.}
\end{quote}

\begin{quote}
The officer scanned my bag.
\end{quote}

Past (I …): \kn{ಸ್ಕ್ಯಾನ್} \kn{ಮಾಡಿದೆ}, \emph{skyān māḍide}.

An English verb plus \kn{ಮಾಡು}, like \kn{ಕ್ಲಿಕ್} \kn{ಮಾಡು}: the thing scanned takes -\kn{ಅನ್ನು}.

\begin{cousinweb}
\textbf{\kn{ಸ್ಕ್ಯಾನ್} \kn{ಮಾಡು}} (\emph{skyān māḍu}) joins a loan from English \emph{scan} and the native verb \kn{ಮಾಡು}, to do.
\end{cousinweb}

\begin{grammarlens}[title={What you've built}]
One more word for getting around.
\end{grammarlens}

\subsection*{Your turn}
\begin{itemize}
\raggedright
  \item \emph{Say it:} \emph{skyān māḍu}
  \item \emph{Say it:} \emph{skyān māḍu}, once more
  \item \emph{Say it:} \emph{adhikāri nanna cīlavannu skyān māḍidaru.} — The officer scanned my bag.
  \item \emph{Recall:} say the Kannada for a terminal, then the Kannada for a boarding pass, then say \emph{skyān māḍu} again
\end{itemize}

\begin{tcolorbox}[breakable,colback=teal!4,colframe=teal!35!black,title={Before you move on}]
What does \kn{ಸ್ಕ್ಯಾನ್} \kn{ಮಾಡು} mean? (\textquotedblleft{}To scan\textquotedblright{}.) Say it once more.
\end{tcolorbox}
\section[\texorpdfstring{\kn{ಪೈಲಟ್} (pailaṭ)}{ಪೈಲಟ್ (pailaṭ)}]{\kn{ಪೈಲಟ್} (pailaṭ) — a pilot}

\label{lesson:KA-C437-pailat}



\begin{quote}
Before the new one: say the Kannada for a boarding pass, then the Kannada for to scan.
\end{quote}

\subsection*{You'll want to know: \kn{ಪೈಲಟ್}}
\textbf{\kn{ಪೈಲಟ್}} — \emph{pailaṭ} — \textquotedblleft{}a pilot\textquotedblright{}.

\begin{quote}
\kn{ನನ್ನ} \kn{ಮಾವ} \kn{ವಿಮಾನದ} \kn{ಪೈಲಟ್} \kn{ಆಗಿದ್ದಾರೆ}.
\end{quote}

\begin{quote}
\emph{nanna māva vimānada pailaṭ āgiddāre.}
\end{quote}

\begin{quote}
My uncle is an aeroplane pilot.
\end{quote}

Plural: \kn{ಪೈಲಟ್ಗಳು}, \emph{pailaṭgaḷu}.

The plural takes -\kn{ಗಳು}: \kn{ಪೈಲಟ್ಗಳು}.

\begin{cousinweb}
\textbf{\kn{ಪೈಲಟ್}} (\emph{pailaṭ}) is a loan from English \emph{pilot}.
\end{cousinweb}

\begin{grammarlens}[title={What you've built}]
One more word for getting around.
\end{grammarlens}

\subsection*{Your turn}
\begin{itemize}
\raggedright
  \item \emph{Say it:} \emph{pailaṭ}
  \item \emph{Say it:} \emph{pailaṭ}, once more
  \item \emph{Say it:} \emph{nanna māva vimānada pailaṭ āgiddāre.} — My uncle is an aeroplane pilot.
  \item \emph{Recall:} say the Kannada for a boarding pass, then the Kannada for to scan, then say \emph{pailaṭ} again
\end{itemize}

\begin{tcolorbox}[breakable,colback=teal!4,colframe=teal!35!black,title={Before you move on}]
What does \kn{ಪೈಲಟ್} mean? (\textquotedblleft{}A pilot\textquotedblright{}.) Say it once more.
\end{tcolorbox}
